\documentclass[10pt,a4paper]{amsart}
\usepackage{amsmath,amssymb,amsthm,graphicx}
% Retained preamble for the Stacks Project chapter snapshot in this directory.
% The chapter's own preamble.tex is a copy of the shared Stacks Project preamble; the block
% below is sliced MECHANICALLY, line for line, from it. Nothing was added, removed or
% redefined; the wrapper class, the xr-hyper/hyperref wiring and the externaldocument list
% are omitted because the assembler supplies the document class.
%
% lines taken from preamble.tex: 3, 5, 9, 11, 12, 156, 157, 158, ... (51 lines total)
%
% package coverage restored from the chapter's own preamble (2026-10-13): xy, xy-2cell

\usepackage[all]{xy}
\xyoption{2cell}
\UseAllTwocells

\newenvironment{reference}{\comment}{\endcomment}
\newenvironment{slogan}{\comment}{\endcomment}
\newenvironment{history}{\comment}{\endcomment}
\theoremstyle{plain}
\newtheorem{theorem}[subsection]{Theorem}
\newtheorem{proposition}[subsection]{Proposition}
\newtheorem{lemma}[subsection]{Lemma}
\theoremstyle{definition}
\newtheorem{definition}[subsection]{Definition}
\newtheorem{example}[subsection]{Example}
\newtheorem{exercise}[subsection]{Exercise}
\newtheorem{situation}[subsection]{Situation}
\theoremstyle{remark}
\newtheorem{remark}[subsection]{Remark}
\newtheorem{remarks}[subsection]{Remarks}
\numberwithin{equation}{subsection}
\def\lim{\mathop{\mathrm{lim}}\nolimits}
\def\colim{\mathop{\mathrm{colim}}\nolimits}
\def\Spec{\mathop{\mathrm{Spec}}}
\def\Hom{\mathop{\mathrm{Hom}}\nolimits}
\def\Ext{\mathop{\mathrm{Ext}}\nolimits}
\def\SheafHom{\mathop{\mathcal{H}\!\mathit{om}}\nolimits}
\def\SheafExt{\mathop{\mathcal{E}\!\mathit{xt}}\nolimits}
\def\Sch{\mathit{Sch}}
\def\Mor{\mathop{\mathrm{Mor}}\nolimits}
\def\Ob{\mathop{\mathrm{Ob}}\nolimits}
\def\Sh{\mathop{\mathit{Sh}}\nolimits}
\def\NL{\mathop{N\!L}\nolimits}
\def\CH{\mathop{\mathrm{CH}}\nolimits}
\def\proetale{{pro\text{-}\acute{e}tale}}
\def\etale{{\acute{e}tale}}
\def\QCoh{\mathit{QCoh}}
\def\Ker{\mathop{\mathrm{Ker}}}
\def\Im{\mathop{\mathrm{Im}}}
\def\Coker{\mathop{\mathrm{Coker}}}
\def\Coim{\mathop{\mathrm{Coim}}}
\DeclareMathSymbol{\boxtimes}{\mathbin}{AMSa}{"02}
\def\QCohstack{\mathcal{QC}\!\mathit{oh}}
\def\Cohstack{\mathcal{C}\!\mathit{oh}}
\def\Spacesstack{\mathcal{S}\!\mathit{paces}}
\def\Quotfunctor{\mathrm{Quot}}
\def\Hilbfunctor{\mathrm{Hilb}}
\def\Curvesstack{\mathcal{C}\!\mathit{urves}}
\def\Polarizedstack{\mathcal{P}\!\mathit{olarized}}
\def\Complexesstack{\mathcal{C}\!\mathit{omplexes}}
\def\Pic{\mathop{\mathrm{Pic}}\nolimits}
\def\Picardstack{\mathcal{P}\!\mathit{ic}}
\def\Picardfunctor{\mathrm{Pic}}
\def\Deformationcategory{\mathcal{D}\!\mathit{ef}}
\usepackage{cleveref}
\crefname{originallemma}{original lemma}{original lemmas}
\Crefname{originallemma}{Original lemma}{Original lemmas}
\makeatletter
\newlabel{lemma-reduce-to-scheme-base}{{0GFT}{source}{}{sourcepointer-1}{}}
\newlabel{lemma-reduce-to-scheme-base@cref}{{[originallemma][0][]0GFT}{[1][0][]source}{}{}{}}
\makeatother
\crefname{originallemma}{original lemma}{original lemmas}
\Crefname{originallemma}{Original lemma}{Original lemmas}
\makeatletter
\newlabel{lemma-reduce-to-scheme-source}{{0GFV}{source}{}{sourcepointer-2}{}}
\newlabel{lemma-reduce-to-scheme-source@cref}{{[originallemma][0][]0GFV}{[1][0][]source}{}{}{}}
\makeatother
\crefname{originallemma}{original lemma}{original lemmas}
\Crefname{originallemma}{Original lemma}{Original lemmas}
\makeatletter
\newlabel{spaces-limits-proposition-there-is-a-scheme-finite-over}{{09YC}{source}{}{sourcepointer-3}{}}
\newlabel{spaces-limits-proposition-there-is-a-scheme-finite-over@cref}{{[originallemma][0][]09YC}{[1][0][]source}{}{}{}}
\makeatother
\crefname{originallemma}{original lemma}{original lemmas}
\Crefname{originallemma}{Original lemma}{Original lemmas}
\makeatletter
\newlabel{spaces-morphisms-lemma-base-change-proper}{{04WP}{source}{}{sourcepointer-4}{}}
\newlabel{spaces-morphisms-lemma-base-change-proper@cref}{{[originallemma][0][]04WP}{[1][0][]source}{}{}{}}
\makeatother
\crefname{originallemma}{original lemma}{original lemmas}
\Crefname{originallemma}{Original lemma}{Original lemmas}
\makeatletter
\newlabel{spaces-morphisms-lemma-base-change-surjective}{{03MH}{source}{}{sourcepointer-5}{}}
\newlabel{spaces-morphisms-lemma-base-change-surjective@cref}{{[originallemma][0][]03MH}{[1][0][]source}{}{}{}}
\makeatother
\crefname{originallemma}{original lemma}{original lemmas}
\Crefname{originallemma}{Original lemma}{Original lemmas}
\makeatletter
\newlabel{spaces-morphisms-lemma-separated-over-separated}{{04ZH}{source}{}{sourcepointer-6}{}}
\newlabel{spaces-morphisms-lemma-separated-over-separated@cref}{{[originallemma][0][]04ZH}{[1][0][]source}{}{}{}}
\makeatother
\crefname{originallemma}{original lemma}{original lemmas}
\Crefname{originallemma}{Original lemma}{Original lemmas}
\makeatletter
\newlabel{spaces-morphisms-section-quasi-compact}{{03HC}{source}{}{sourcepointer-7}{}}
\newlabel{spaces-morphisms-section-quasi-compact@cref}{{[originallemma][0][]03HC}{[1][0][]source}{}{}{}}
\makeatother
\title{Etale sheaves satisfy descent along surjective integral, proper, or fppf-flat morphisms}
\author{The Stacks Project Authors}
\date{stacks-project:spaces-pushouts.tex:local-snapshot-20261012; GNU FDL 1.2 or later}
\begin{document}
\maketitle
\noindent\textit{Editorial note.} Statement span is lines 816-828 of sources/stacks-spaces-pushouts/pinned-original-source.tex; proof span is lines 830-863, the complete author proof (32 body lines).\par\medskip
\noindent\textit{Editorial note.} Navigation stubs: 7 cited-but-not-included labels are supplied as non-load-bearing pointers with their pinned tags where the snapshot tag index has them.\par\medskip
\noindent\textit{Editorial note.} extra\_packages is empty on purpose: the retained preamble already loads xy with options and a second declaration clashes in the assembler wrapper.\par\medskip
\noindent\textit{Editorial note.} Text is retained verbatim, including the author's own typos and placeholders.\par\medskip
\noindent\textit{Editorial note.} Queue index 19 of queue\_spaces-pushouts.json (pb=1450, stmt 816-828, proof 830-863).\par\medskip
\noindent\textit{Editorial note.} spaces-pushouts wave-1: fresh adjudication accept (2026-10-14); six-check dispatched pool (§2 row appended 02:4x).\par\medskip
\noindent\textit{Allowed standard prerequisites (editorial).}\begin{itemize}
\item Explicit prerequisite boundary. The cited inputs are prerequisites; the argument in the retained span is not replaced by any cited result.
\end{itemize}\medskip
\ifdefined\hypertarget\hypertarget{sourcepointer-1}{}\fi
\noindent\textit{Source reference.} Original lemma 0GFT: sources/stacks-spaces-pushouts/pinned-original-source.tex, source lines 729--752. Cited by the retained statement or proof but not part of the collected spans; supplied only as a navigation pointer. The argument does not depend on its content.\par\medskip
\ifdefined\hypertarget\hypertarget{sourcepointer-2}{}\fi
\noindent\textit{Source reference.} Original lemma 0GFV: sources/stacks-spaces-pushouts/pinned-original-source.tex, source lines 792--801. Cited by the retained statement or proof but not part of the collected spans; supplied only as a navigation pointer. The argument does not depend on its content.\par\medskip
\ifdefined\hypertarget\hypertarget{sourcepointer-3}{}\fi
\noindent\textit{Source reference.} Original lemma 09YC: sources/stacks-spaces-pushouts/pinned-original-source.tex, source lines 863--863. Cited by the retained statement or proof but not part of the collected spans; supplied only as a navigation pointer. The argument does not depend on its content.\par\medskip
\ifdefined\hypertarget\hypertarget{sourcepointer-4}{}\fi
\noindent\textit{Source reference.} Original lemma 04WP: sources/stacks-spaces-pushouts/pinned-original-source.tex, source lines 863--863. Cited by the retained statement or proof but not part of the collected spans; supplied only as a navigation pointer. The argument does not depend on its content.\par\medskip
\ifdefined\hypertarget\hypertarget{sourcepointer-5}{}\fi
\noindent\textit{Source reference.} Original lemma 03MH: sources/stacks-spaces-pushouts/pinned-original-source.tex, source lines 863--863. Cited by the retained statement or proof but not part of the collected spans; supplied only as a navigation pointer. The argument does not depend on its content.\par\medskip
\ifdefined\hypertarget\hypertarget{sourcepointer-6}{}\fi
\noindent\textit{Source reference.} Original lemma 04ZH: sources/stacks-spaces-pushouts/pinned-original-source.tex, source lines 863--863. Cited by the retained statement or proof but not part of the collected spans; supplied only as a navigation pointer. The argument does not depend on its content.\par\medskip
\ifdefined\hypertarget\hypertarget{sourcepointer-7}{}\fi
\noindent\textit{Source reference.} Original lemma 03HC: sources/stacks-spaces-pushouts/pinned-original-source.tex, source lines 863--863. Cited by the retained statement or proof but not part of the collected spans; supplied only as a navigation pointer. The argument does not depend on its content.\par\medskip
\begin{lemma}
\label{lemma-glue-etale-sheaf-proper-surjective}
Let $S$ be a scheme. Let $f : X \to Y$ be a morphism of
algebraic spaces over $S$ which has one of the following
properties: surjective and integral, surjective and proper, or
surjective and flat and locally of finite presentation. Then the functor
$$
\Sh(Y_\etale)
\longrightarrow
\text{descent data for \'etale sheaves wrt }\{X \to Y\}
$$
is an equivalence of categories.
\end{lemma}

\begin{proof}
Observe that the base change of a proper surjective morphism is
proper and surjective, see
Morphisms of Spaces, Lemmas \ref{spaces-morphisms-lemma-base-change-proper}
and \ref{spaces-morphisms-lemma-base-change-surjective}.
Hence by Lemma \ref{lemma-reduce-to-scheme-base}
we may work \'etale locally on $Y$. Hence
we reduce to $Y$ being an affine scheme; some details omitted.

\medskip\noindent
Assume $Y$ is affine. By Lemma \ref{lemma-reduce-to-scheme-source}
it suffices to find a morphism $X' \to X$ where $X'$ is a scheme such
that $X' \to Y$ is surjective and integral, surjective and proper, or
surjective and flat and locally of finite presentation.

\medskip\noindent
In case $X \to Y$ is integral and surjective, we can take $X = X'$
as an integral morphism is representable.

\medskip\noindent
If $f$ is proper and surjective, then the algebraic space
$X$ is quasi-compact and separated, see
Morphisms of Spaces, Section \ref{spaces-morphisms-section-quasi-compact} and
Lemma \ref{spaces-morphisms-lemma-separated-over-separated}.
Choose a scheme $X'$ and a surjective finite morphism $X' \to X$, see
Limits of Spaces, Proposition
\ref{spaces-limits-proposition-there-is-a-scheme-finite-over}.
Then $X' \to Y$ is surjective and proper.

\medskip\noindent
Finally, if $X \to Y$ is surjective and flat and locally of finite
presentation then we can take an affine \'etale covering $\{U_i \to X\}$
and set $X'$ equal to the disjoint $\coprod U_i$.
\end{proof}

\end{document}
